% Ejemplo del estilo de informe UD. Compilar desde plantillas-latex/:  make ejemplos
\documentclass[11pt,letterpaper]{article}
\usepackage{ud-informe}

\renewcommand{\UDtitulo}{Título del informe}
\renewcommand{\UDsubtitulo}{Subtítulo que explica el alcance}
\renewcommand{\UDtipo}{Taller 1}
\renewcommand{\UDasignatura}{Asignatura}
\renewcommand{\UDdocente}{Nombre del docente}
\renewcommand{\UDestudiantes}{Juan Felipe Rodríguez Galindo -- 20261595004}
\renewcommand{\UDencabezado}{Taller 1 -- Asignatura}

\begin{document}
\UDportada
\tableofcontents
\newpage

\section{Introducción}

Texto del informe en Times New Roman; los títulos van en Cambria con el número en rojo institucional.

\begin{table}[!htbp]
\centering
\caption{Una tabla de ejemplo.}
\begin{tabularx}{\linewidth}{l L r}
\toprule
Columna & Descripción & Valor \\
\midrule
A & Texto que se ajusta al ancho disponible & 1,0 \\
\bottomrule
\end{tabularx}
\end{table}

\begin{figure}[!htbp]
\centering
\begin{tikzpicture}[node distance=0.6cm]
  \node[caja] (a) {Datos}; \node[caja, right=of a] (b) {Modelo}; \node[hecho, right=of b] (c) {Resultado};
  \draw[flecha] (a) -- (b); \draw[flecha] (b) -- (c);
\end{tikzpicture}
\caption{Un diagrama con los estilos \texttt{caja}, \texttt{hecho} y \texttt{flecha}.}
\end{figure}

\end{document}
